\subsection{Codimension of a closed subset}

DEFINITION 3. --- Let X be a topological space.

\begin{enumerate}
\def\labelenumi{\alph{enumi})}
\item
  If Y is a closed irreducible subset of X, one calls the codimension of Y in X the supremum in \(\overline{\mathbf{R}}\) of the lengths of the chains of irreducible closed subsets of X for which Y is the smallest element.
\item
  If Y is a closed subset of X, we call the codimension of Y in X, and denote it by \(\text{codim}(Y, X)\) the infimum in \(\overline{\mathbf{R}}\) of the codimensions, in X, of the irreducible components of Y.
\end{enumerate}

Remarks. --- 1) The codimension of a closed subset Y of X is therefore the infimum of the codimensions of the irreducible closed subsets of Y. We have \(\mathrm{codim}(\emptyset, \mathrm{X}) = +\infty\) and, if X is nonempty, \(\mathrm{codim}(\mathrm{X}, \mathrm{X}) = 0\) . Every nonempty closed subset of X contains an irreducible closed subset (II, § 4, no. 1, prop. 5); hence the codimension in X of a closed subset Y is always a nonnegative integer, or \(+\infty\) . It is zero if and only if Y contains an irreducible component of X.

\begin{enumerate}
\def\labelenumi{\arabic{enumi})}
\setcounter{enumi}{1}
\tightlist
\item
  If Y is a nonempty closed subset of X, then codim(Y, X) ≤ dim(X). Moreover, dim(X) = sup codim(Y, X), where Y ranges over the set of irreducible closed subsets.
\end{enumerate}

If Y and Y' are two closed subsets of X such that Y' ⊂ Y, then codim(Y, X) ≤ codim(Y', X).

\begin{enumerate}
\def\labelenumi{\arabic{enumi})}
\setcounter{enumi}{2}
\tightlist
\item
  Let Y be a closed subset of the topological space X and \((\mathrm{X}_{\alpha})_{\alpha\in\mathrm{A}}\) (resp. \((\mathrm{Y}_{\beta})_{\beta\in\mathrm{B}}\) ) the family of irreducible components of X (resp. of Y). For every \(\beta\in B\) , denote by \(\mathrm{A}(\beta)\) the set of \(\alpha\in A\) such that \(Y_{\beta}\subset X_{\alpha}\) . Since every irreducible subset of X is contained in one of the \(X_{\alpha}\) (II, § 4, no. 1, prop. 5), it follows from Def. 3 that we have:
\end{enumerate}

\[
\operatorname{codim} (Y, X) = \inf _ {\beta \in B} \sup _ {\alpha \in A (\beta)} \operatorname{codim} (Y _ {\beta}, X _ {\alpha}).
\]

\begin{enumerate}
\def\labelenumi{\arabic{enumi})}
\setcounter{enumi}{3}
\tightlist
\item
  Let \((\mathbf{Y}_i)_{i\in I}\) be a finite family of closed subsets of X and \(\mathbf{Y} = \bigcup_{i\in I}\mathbf{Y}_i\) ; one has
\end{enumerate}

\[
\operatorname{codim} (\mathrm{Y}, \mathrm{X}) = \inf _ {i \in \mathrm{I}} \operatorname{codim} (\mathrm{Y} _ {i}, \mathrm{X}).
\]

Indeed, every irreducible component of Y is contained in one of the \(Y_{i}\) .

PROPOSITION 3. --- Let X be a topological space.

\begin{enumerate}
\def\labelenumi{\alph{enumi})}
\tightlist
\item
  For every nonempty closed subset Y of X, we have
\end{enumerate}

\[
\dim (\mathrm{Y}) + \operatorname{codim} (\mathrm{Y}, \mathrm{X}) \leqslant \dim (\mathrm{X}).
\]

\begin{enumerate}
\def\labelenumi{\alph{enumi})}
\setcounter{enumi}{1}
\tightlist
\item
  If Y, Z, and T are closed subsets of X such that Y ⊂ Z ⊂ T, then
\end{enumerate}

\[
\operatorname{codim} (\mathrm{Y}, \mathrm{Z}) + \operatorname{codim} (\mathrm{Z}, \mathrm{T}) \leqslant \operatorname{codim} (\mathrm{Y}, \mathrm{T}).
\]

It suffices to prove the assertion \(a)\) in the case where \(\dim(X)\) is finite. In this case, \(\dim(Y)\) and \(\operatorname{codim}(Y, X)\) are finite. There exists a chain \(Y_0 \subset \ldots \subset Y_n\) of closed irreducible subsets of \(Y\), of length \(n = \dim(Y)\) and a chain \(Y_n \subset \ldots \subset Y_{n+p}\) of closed irreducible subsets of \(X\), of length \(p \geqslant \operatorname{codim}(Y, X)\). It follows that \(\dim(X) \geqslant n + p\), whence . To establish \(a)\)\(b)\), one may assume \(Y\) irreducible. Since we have \(\operatorname{codim}(Y, Z) \leqslant \operatorname{codim}(Y, T)\), the inequality is proved if \(\operatorname{codim}(Y, Z) = +\infty\). Otherwise, let \(Z_0\) be an irreducible component of \(Z\) containing \(Y\) and such that \(\operatorname{codim}(Y, Z) = \operatorname{codim}(Y, Z_0)\). We have \(\operatorname{codim}(Z, T) \leqslant \operatorname{codim}(Z_0, T)\), and one sees, as above, that \(\operatorname{codim}(Y, Z_0) + \operatorname{codim}(Z_0, T) \leqslant \operatorname{codim}(Y, T)\), whence \(b)\).

DEFINITION 4. --- A topological space X is said to be catenary if, for every pair (Y, Z) of irreducible closed subsets of X such that Y ⊂ Z, every saturated chain of irreducible closed subsets with endpoints Y and Z has length codim(Y, Z).

It is equivalent to say that, for every pair (Y, Z) of irreducible closed subsets of X such that codim(Y, Z) is finite, all saturated chains with endpoints Y and Z have the same length, and that, for every pair (Y, Z) such that codim(Y, Z) = +∞, there exists no saturated chain with endpoints Y and Z.

Every closed subspace of a catenary space is catenary. For a space to be catenary, it is necessary and sufficient that its irreducible components be catenary.

PROPOSITION 4. --- Let X be a topological space. For X to be catenary, it is necessary and sufficient that, for every triple (Y, Z, T) of irreducible closed subsets of X such that Y ⊂ Z ⊂ T, one has:

\[
\operatorname{codim} (\mathrm{Y}, \mathrm{T}) = \operatorname{codim} (\mathrm{Y}, \mathrm{Z}) + \operatorname{codim} (\mathrm{Z}, \mathrm{T}).
\]

Suppose X is catenary. In view of prop. 3, b), it suffices to prove the relation when codim(Y, Z) and codim(Z, T) are finite. By concatenating a saturated chain of irreducible closed subsets with endpoints Y and Z, of length codim(Y, Z), and a saturated chain of irreducible closed subsets with endpoints Z and T, of length codim(Z, T), one obtains a saturated chain with endpoints Y and T, of length codim(Y, Z) + codim(Z, T). But, since X is catenary, this length is necessarily equal to codim(Y, T).

Conversely, suppose we have \(\operatorname{codim}(\mathbf{Y},\mathbf{T})=\operatorname{codim}(\mathbf{Y},\mathbf{Z})+\operatorname{codim}(\mathbf{Z},\mathbf{T})\) for all irreducible closed subsets  of  such that \(\mathbf{Y}\subset\mathbf{Z}\subset\mathbf{T}\), and let us prove that X is catenary. To do this, we prove by induction on the integer \(n\geqslant0\) that, for every saturated chain \(Z_{0}\subset\ldots\subset Z_{n}\) of closed irreducible subsets of X, we have \(\operatorname{codim}(Z_{0},Z_{n})=n\). If \(n=0\), it is clear. Let \(n>0\), and suppose the property is satisfied for chains of length \(\leqslant n-1\). If \(Z_{0}\subset\ldots\subset Z_{n}\) is a saturated chain of length \(n\), then \(Z_{0}\subset\ldots\subset Z_{n-1}\) is a saturated chain of length \(n-1\), hence \(\operatorname{codim}(Z_{0},Z_{n-1})=n-1\). In view of the hypothesis made on X, we have \(\operatorname{codim}(Z_{0},Z_{n})=\operatorname{codim}(Z_{0},Z_{n-1})+\operatorname{codim}(Z_{n-1},Z_{n})=(n-1)+1=n\).

COROLLARY. — Let X be an irreducible topological space of finite dimension. For X to be catenary, it is necessary and sufficient that, for every pair (Y, Z) of irreducible closed subsets of X such that Y ⊂ Z, one has codim(Y, X) = codim(Y, Z) + codim(Z, X).

The condition is necessary by prop. 4. Conversely, suppose it is satisfied, and denote by \(c(Z)\) the integer \(\operatorname{codim}(Z,X)\) for every irreducible closed subset \(Z\) of \(X\) . If \(Y,Z,T\) are three irreducible closed subsets of \(X\) such that \(Y\subset Z\subset T\) , one has

\[
\begin{array}{r l} \operatorname{codim} (\mathrm{Y}, \mathrm{Z}) + \operatorname{codim} (\mathrm{Z}, \mathrm{T}) & = (c (\mathrm{Y}) - c (\mathrm{Z})) + (c (\mathrm{Z}) - c (\mathrm{T})) \\ & = c (\mathrm{Y}) - c (\mathrm{T}) \\ & = \operatorname{codim} (\mathrm{Y}, \mathrm{T}), \end{array}
\]

and X is catenary by prop. 4.

PROPOSITION 5. — Let X be a topological space of finite dimension. Suppose that all maximal chains of irreducible closed subsets of X have the same length. Then X is catenary; for every irreducible closed subset Z of X, one has

\[
\operatorname{codim} (Z, X) = \dim (X) - \dim (Z);
\]

for every pair (Y, Z) of irreducible closed subsets of X such that Y ⊂ Z, one has

\[
\operatorname{codim} (\mathrm{Y}, \mathrm{Z}) = \dim (\mathrm{Z}) - \dim (\mathrm{Y}).
\]

Let Y and Z be two irreducible closed subsets of X such that \(Y \subset Z\) . Let \(Y_{0} \subset \ldots \subset Y_{p}\) be a chain such that \(Y_{p} = Y\) and \(p = \dim(Y)\) , \(Z_{0} \subset \ldots \subset Z_{q}\) be a chain such that \(Z_{0} = Z\) and \(q = \text{codim}(Z, X)\) . For every saturated chain \(T_{0} \subset \ldots \subset T_{r}\) such that \(T_{0} = Y\) and \(T_{r} = Z\) , the chain

\[
\mathrm{Y} _ {0} \subset \dots \subset \mathrm{Y} _ {p - 1} \subset \mathrm{T} _ {0} \subset \dots \subset \mathrm{T} _ {r} \subset \mathrm{Z} _ {1} \dots \subset \mathrm{Z} _ {q}
\]

is maximal, and of length \(p + q + r\) ; by the hypothesis made on X, one therefore has \(p + q + r = \dim(X)\) . Let \(r = \dim(X) - \dim(Y) - \text{codim}(Z, X)\) . It follows that X is catenary and that, for Y and Z as above, one has

\[
\dim (\mathrm{Y}) + \operatorname{codim} (\mathrm{Y}, \mathrm{Z}) = \dim (\mathrm{X}) - \operatorname{codim} (\mathrm{Z}, \mathrm{X}).
\]

Taking Y = Z, one sees that the second member is equal to dim(Z), whence the proposition.

